\section{Rappels de première année}
\subsection{Définition}

\begin{definition}{Espace préhilbertien réel}{}
    On appelle espace préhilbertien réel un espace vectoriel réel $E$
    sur $\R$ muni d'un produit scalaire, c'est à dire d'une forme 
    bilinéaire symétrique définie positive.
\end{definition}

\begin{proposition}{}{}
    Soit $(E, \langle \cdot, \cdot \rangle)$ un espace préhilbertien réel. Alors
    \begin{itemize}
        \item Inégalité de Cauchy-Schwarz : $\forall x, y \in E, |\langle x, y \rangle| \leqslant \|x\|_2 \|y\|_2$ avec 
        égalité ssi $x$ et $y$ sont liés.
        \item La norme $\|x\|_2 = \sqrt{\langle x, x \rangle}$ est une norme appellée norme euclidienne.
        \item $\forall u, v \in E$, $\|u + v\|_2^2 = \|u\|_2^2 + 2\langle u, v \rangle + \|v\|_2^2$.
        \item $\forall u, v \in E$, $\|u - v\|_2^2 = \|u\|_2^2 - 2\langle u, v \rangle + \|v\|_2^2$.
        \item Identité de polarisation : $\forall u, v \in E$, $\langle u, v \rangle = \frac{1}{4}(\|u + v\|_2^2 - \|u - v\|_2^2)$.
        \item Identité du parallélogramme : $\forall u, v \in E$, $\|u + v\|_2^2 + \|u - v\|_2^2 = 2\|u\|_2^2 + 2\|v\|_2^2$.
    \end{itemize}
\end{proposition}

\subsection{Orthogonalité}

\begin{definition}{Orthogonalité}{}
    Dans $(E, \langle \cdot, \cdot \rangle)$ un esapce préhilbertien réel, on dit que 
    $u$ et $v$ sont orthogonaux et on note $u \perp v$ ssi $\langle u, v \rangle = 0$.
\end{definition}

\begin{definition}{Orthogonalité de deux ensembles}{}
    Soient $A, B \subset E$. Alors $A \perp B$ ssi $\forall a \in A, \, \forall b \in B$,
    $a \perp b$.
\end{definition}

\begin{definition}{Orthogonal d'un ensemble}{}
    Soit $A \subset E$. $A^\perp = \{x \in E | \forall a \in A, x \perp a\}$. 
    Ainsi, $A \perp A^\perp$. 

    Alors $A^\perp$ est un sev de $E$ et $ A \subset (A^\perp)^\perp$.
\end{definition}

\begin{theoreme}{Projection orthogonale sur un sous-espace vectoriel}{}
    Soit $(E, \langle \cdot, \cdot \rangle)$ un espace préhilbertien réel, 
    et soit $F \subset E$ un sous-espace vectoriel de dimension finie. Alors 
    $E = F \oplus F^\perp$. Alors $F^\perp$ est appelé le supplémentaire 
    orthogonal de $F$. 

    De plus en notant $p$ la projection sur $F$ parallèlement à $F^\perp$, $p$ est 
    appelée projection orthogonale sur $F$ et on a : 
    $\forall x \in E, \, x = p(x) + x - p(x)$ et $\| x - p(x) \|_2 = \min_{u \in F} \| x - u \|_2$.
    $p(x)$ est l'unique vecteur de $F$ vérifiant l'égalité. 

    Si $(e_i)$ forment une base orthonormale de $F$, on a : 
    $$ p(x) = \sum_{i = 1}^k \langle x, e_i \rangle e_i$$
\end{theoreme}

\begin{proof}
    $\blacktriangleright$ Soit $x \in F \cap F^\perp$, alors $\langle x, x \rangle = 0$ et donc $x = 0$. 
    Donc la somme est directe. 

    \avspace $\blacktriangleright$ Soit $x \in E$, soit $(e_i)$ une base de $F$. 

    Posons $y = \ds \sum_{i = 1}^k \langle x, e_i \rangle e_i$ et $z = x - y$. 

    On a $y \in F$, montrons que $z \in F^\perp$. 

    Soit $j \in \ninter 1 k$. 

    \begin{align}
        \langle z, e_j \rangle & = \langle x - y, e_j \rangle \\
        & = \langle x, e_j \rangle - \sum_{i = 1}^k \langle x, e_i \rangle \langle e_i, e_j \rangle \\
        & = 0
    \end{align}

    Donc $z \perp e_j$ et ainsi $z \in F^\perp$.

    Donc $E \subset F \oplus F^\perp$.

    Donc $E = F \oplus F^\perp$ et $\forall x \in E$, $p(x) = y$. 

    \avspace $\blacktriangleright$ Soit $u \in F$, alors :
    \begin{align}
        \| x - u \|_2^2 & = \| (x - p(x)) + (p(x) - u) \|_2^2 \\
        & = \| x - p(x) \|_2^2 + 2\langle x - p(x), p(x) - u \rangle + \| p(x) - u \|_2^2 \\
        & \geqslant || x - p(x) \|_2^2
    \end{align}

    avec égalité ssi $p(x) = u$. 

    $p(x)$ est donc l'unique vecteur de $F$ minimisant la distance à $x$.

    \avspace $\blacktriangleright$ On a $F \subset (F^\perp)^\perp$. Soit $x \in (F^\perp)^\perp$,
    et décomposons le en $x = y + z$ avec $y \in F$ et $z \in F^\perp$.

    On a donc $0 = \langle x, z \rangle$ or $\langle x, z \rangle = \langle y, z \rangle + \langle z, z \rangle = \|z\|_2^2$.
    Donc $z = 0$ et $x = y \in F$. Donc $(F^\perp)^\perp \subset F$ et ainsi $F = (F^\perp)^\perp$.
\end{proof}

\begin{exemple}
    Chercher un min ou délire du genre en utilisant projection orthogonale (sachant que possible
    d'essayer avec calcul diff mais super long et pénible).

    $E = \mathscr C([0, 1], \R)$ muni du produit scalaire $\langle f, g \rangle = \ds \int_0^1 f(t)g(t) dt$.

    $F = \Vect (u_0, v_0)$ avec $u_0(t) = t$ et $v_0(t) = 1$ et $w_0(t) = e^t$. 

    Alors 
    \begin{align}
        f(a, b) & = \int_0^1 (e^t -at - b)^2 dt \\
        & = \| w_0 - au_0 - bv_0 \|_2^2 
    \end{align}
    distance au carré entre $a u_0 + b v_0 \in F$ quelconque. 

    Donc $f(a, b)$ minimal ssi $au_0 + bv_0 = p(w_0)$ avec $p$ projection orthogonale sur 
    $F$. 

    On a $p(w_0) = a u_0 + b v_0 \in F$

    $\begin{cases}
        w_0 - p(w_0) \perp u_0 \\
        w_0 - p(w_0) \perp v_0
    \end{cases}$ 
    car $w_0 - p(w_0) \in F^\perp$.

    Donc $\begin{cases}
        \langle w_0, u_0 \rangle = \langle p(w_0), u_0 \rangle \\
        \langle w_0, v_0 \rangle = \langle p(w_0), v_0 \rangle
    \end{cases}$
    donc $\begin{cases}
        a \|u_0\|_2^2 + b \langle v_0, u_0 \rangle = \langle w_0, u_0 \rangle \\
        a \langle u_0, v_0 \rangle + b \|v_0\|_2^2 = \langle w_0, v_0 \rangle
    \end{cases}$

    Donc $\ds \| u_0 \|^2 = \int_0^1 t^2 dt = \frac{1}{3}$, $\| v_0 \|^2 = 1$, $\langle u_0, v_0 \rangle = \int_0^1 t dt = \frac{1}{2}$, $\langle w_0, u_0 \rangle = \int_0^1 t e^t dt$ et $\langle w_0, v_0 \rangle = \int_0^1 e^t dt$.
    Donc $\ds a \frac{1}{3} + b \frac{1}{2} = \int_0^1 t e^t dt$ et $a \frac{1}{2} + b = \int_0^1 e^t dt$.

\end{exemple}

\begin{exemple}
    On a $(x_i, y_i) \in \mathbb R^2$ et on cherche $a, b$ tels que 
    $\ds \sum_{i = 1}^n (y_i - ax_i - b)^2$ soit minimal.

    Posons $\fonction f {\R^2}{\R}{(a, b)}{\ds \sum_{i = 1}^n (y_i - ax_i - b)^2}$.

    Posons $Y = \begin{pmatrix}
y_1 \\
\vdots \\
y_n
\end{pmatrix}$, $X = \begin{pmatrix}
    x_1 \\
    \vdots \\
    x_n
    \end{pmatrix}$ et $U = \begin{pmatrix}
    1 \\
    \vdots \\
    1
    \end{pmatrix}$ en munissant $\R^n$ du produit scalaire canonique, 
    $f(a, b) = \| Y - (a X + bU) \|^2$ 

    $Y \in \R^n$ et $(X, U)$ famille libre de $\R^n$ donc $F = \Vect(X, U)$ est un plan. 

    Le minimum de $f$ est donc atteint en $(a, b)$ pour lequel $aX + bU$ est le projeté
    orthogonal de $Y$ sur $F$. 

    Il est caractérisé par $Y = \underbrace{aX + bU}_{\in F} + \underbrace{(Y - aX - bU)}_{\in F^\perp}$, 
    c'est à dire que 
    $\begin{cases}
        \langle Y - (aX + bU), X \rangle = 0 \\
        \langle Y - (aX + bU), U \rangle = 0
    \end{cases}$
    
    c'est à dire que $\begin{cases}
        a \|X\|^2 + b \langle U, X \rangle = \langle Y, X \rangle \\
        a \langle X, U \rangle + b \|U\|^2 = \langle Y, U \rangle
    \end{cases}$

    Notons $m_X = \frac 1 n \ds \sum_{i = 1}^n x_i$ et $m_Y = \frac 1 n \ds \sum_{i = 1}^n y_i$.

    $v_X = \frac 1 n \ds \sum_{i = 1 }^{n } x_i^2 - m_X^2$

    Et donc $\operatorname{Cov} (X, Y) = \frac 1 n \ds \sum_{i = 1}^n x_i y_i - m_X m_Y$.

    $\langle U, X \rangle = \ds \sum_{i = 1}^{n } x_i = n m_x$

    $\| x \|^2 = \ds \sum_{i = 1}^{n } x_i^2 = nv_x + nm_x^2$

    $\langle Y, X \rangle = \ds \sum_{i = 1}^n x_i y_i = n \operatorname{Cov}(X, Y) + nm_X m_Y$

    $\langle Y, U \rangle = \ds \sum_{i = 1}^n y_i = n m_Y$

    $(S) \Leftrightarrow \ds \begin{cases}
a (nv_x + nm_x^2) + b n m_x = n \operatorname{Cov}(X, Y) + nm_X m_Y \\
a n m_x + b n = n m_Y
    \end{cases} \Leftrightarrow \begin{cases}
        am_x + b = m_Y \\
        a v_x = \operatorname{Cov}(X, Y)
    \end{cases}$

    $v_x \neq 0$ car les $x_i$ ne sont pas tous égaux.

    Donc $(S) \Leftrightarrow \begin{cases}
        a = \frac{\operatorname{Cov}(X, Y)}{v_x} \\
        b = m_Y - a m_x
    \end{cases}$
\end{exemple}

\begin{exemple}
    $\ds E = \{ f : \R \to \R \text{ cpm }, 2 \pi \text{ périodique } | \forall x \in \R, f(x) = \frac 1 2 \left( \lim_{x^-} f + \lim_{x^+} f \right) \}$

    On pose $\langle f, g \rangle = \frac{1}{\pi} \int_0^{2\pi} f(t) g(t) dt$.

    Supposons que $\langle f, f \rangle = 0$, alors $\int_{0}^{2 \pi} f(t)^2 dt = 0$ 

    Soit $(c_i)_{0 \leqslant i \leqslant n}$ subdivision de $[0, 2\pi]$ adaptée à $f$ qui est continue par morceaux. 

    $\forall i \in \ninter 1 n$, $f_{|]c_{i - 1}, c_i[} $ est continue et prolongeable par continuité en 
    $c_{i - 1}$ et $c_i$.

    Notons $g_i$ ce prolongement par continuité. 

    Par Chasles, on a $0 = \ds \int_{0}^{2 \pi} f^2 = \sum_{i = 1 }^{n } \int_{c_{i - 1}}^{c_i }f^2(t) \dt$.

    Donc $\forall i \in \ninter 1 n$, on a : $\ds \int_{c_{i - 1}}^{c_i } f^2 (t) \ds = 0$

    Donc $\ds \int_{c_{i - 1}}^{c_i } g_i^2 (t) \ds = 0$ car $f$ et $g_i$ coïncident sauf en deux points. 

    Or $g_i^2 \geqslant 0$ continue, donc $\forall t \in [c_{i - 1}, c_i], g_i(t) = 0$ 

    Donc $\forall t \in ]c_{i - 1}, c_i[, f(t) = 0$.

    De plus par hypothèse, $f(c_i) = \frac 1 2 \left( \lim_{t \to c_i^-} f(t) + \lim_{t \to c_i^+} f(t) \right) = 0$.

    Donc $\forall x \in [0, 2 \pi]$, $f(x) = 0$. Donc $f = 0$ par périodicité. 

    Les autres propriétés sont vérifiées facilement.

    Posons $c_0(t) = \frac 1 {\sqrt 2}$, alors $c_0 \in E$ 
    et $\forall n \in \N^*$, on a $c_n(t) = \cos(nt)$ et $s_n(t) = \sin(nt)$ qui est une famille orthonormale de $E$.

    Posons $E_n = \Vect(c_0, c_1, s_1, \dots, c_n, s_n)$ famille orthonormale de $E$ donc libre. 

    Donc base orthonormale de $E$.

    Soit $f \in E$. Posons $p_n(f) = $ projecteur orthogonal de $f$ sur $E_n$.

    $p_n(f) = \langle c_0, f \rangle c_0 + \ds \sum_{k = 1}^n \langle c_k, f \rangle c_k + \langle s_k, f \rangle s_k$.

    $\langle c_0, f \rangle = \left( \frac 1 \pi \int_{0}^{2 \pi} \frac{f(t)}{\sqrt 2} \dt \right) \frac 1 {\sqrt 2} = \frac{1}{2 \pi} \int_{0}^{2 \pi} f(t) dt$.

    $\langle c_k, f \rangle = \frac 1 \pi \ds \int_{0}^{2 \pi} f(t) \cos(kt) \dt = a_k$ 

    $\langle s_k, f \rangle = \frac 1 \pi \ds \int_{0}^{2 \pi} f(t) \sin(kt) \dt = b_k$. coefs de Fourrier trigonométriques de $f$. 

    $p_n(f)$ est la somme partielle d'une série de fonctions appelée série de Fourier associée 
    à $f$. On montre que $p_n(f) \xrightarrow[n \to + \infty]{\| \cdot \|_2} f$. 

    Si $f$ est $\mc 1$ par morceaux et $f \in E$, $p_n(f) \xrightarrow[n \to + \infty]{CS} f$.

    Si $f$ est continue et $\mc 1$ par morceaux, $p_n(f) \xrightarrow[n \to + \infty]{CU} f$.

    Petits dessins pour la fin 
\end{exemple}

